\section {Formal verification and model checking}
Formal Methods refers to mathematically rigorous techniques and
tools for the specification, design and verification of software and
hardware systems.\\ 
Expanding this definition Formal methods 
can be used for formal description of the system (specification) 
on some level of abstraction, program synthesis (design) according to specification
and verification of a system to prove that model of system under consideration 
satisfies some properties.


These methods and techniques include:
\begin{enumerate}

    \item SAT solving - operate on propositional logic 

    \item SMT solving (and tools like Z3) extend this to first-order logic,
        integrating theories such as arithmetic or arrays.

    \item Relational model finding (e.g., Alloy) 

    \item Model checking (e.g., NuSMV) exhaustively explore state-space against CTL/LTL specifications.

    \item Interactive theorem provers (e.g., Coq/Rocq, Isabelle/HOL) 
\end{enumerate}


\subsection {SAT solvers}

Propositional logic - consider  propositions and relations between them and constructions of argument based on them.

Propositions - facts or statements about the world  that can be true or false. Also defined as atomic propositions or Atoms.
Some propositions can be described in natural language as atomic or indecomposable declarative sentences)

Atomic propositions can be connected to each other with logical connectives and compose propositional formulas.
The most common connective relations are:
\begin{itemize}
    \item Negation
    \item Conjunction
    \item Disunction
    \item Implication
\end{itemize}

Based on these we can construct a system of reasoning: from some set of formulas we can infer (prove) another formula that can be true or false:
\[ \phi_1, \phi_2, \phi3, ... \phi_n \vdash \psi \]
Such formulas called sequents, where on the left side there is a set of formulas called premises and on the right side - a formula called conclusion


Another relation between logical formulas  is semantic entailment:
\[ \phi_1, \phi_2, \phi3, ... \phi_n \models \psi \]
it means that if propositional logic formulas $\phi_1 .. \phi_n$ evaluate to True, then formula $\psi$ also evaluates to True.

The formula $\phi$ is valid if for all valuations of atomic propositions in formula it valuates to True:
\[ \models \phi \]


SAT solver is a program which aims to solve the Boolean satisfiability problem 
(SAT), where the Boolean satisfiability problem sometimes called 
propositional satisfiability asks whether there 
exists an interpretation that satisfies a given Boolean formula. 

Given formula  $\phi$ is \textbf{satisfiable} if it has a valuation in which it valuates to True.:

In other words, it asks whether the formula's variables can be 
consistently replaced by the values TRUE or FALSE to make the formula 
evaluate to TRUE. If this is the case, the formula is called satisfiable, 
else unsatisfiable.\cite{Wik1}
SAT problem is NP-complete, as a result, only algorithms with exponential 
worst-case complexity are known. In spite of this, efficient and 
scalable algorithms for SAT were developed during the 2000s, which have 
contributed to dramatic advances in the ability to automatically 
solve problem instances involving tens of thousands of variables and 
millions of constraints.%[https:en.wikipedia.org/wiki/SAT_solver -1.]

SAT solvers usually begin by converting a problem to conjunctive normal form (CNF).
\textbf{Conjunctive normal form (CNF)} - as a conjunction of clauses, where each clause is a disjunction of literals.
Literal is an atomic proposition or a negated atomic proposition.
Conjunctive normal form allows easy check of validity of a formula which otherwise take times exponential to number of atoms.








\subsection {SMT}
Propositional logic is limited to express relations between logical components other than equivalent to natural language: \emph{not, and, or, if..then} and 
declarative sentences that can only be True or False. 

In order to be able to operate with finer logical structures we need to enrich the language.
The necessity to solve this problem lead to design of Predicate logic or first-order logic.\\
In First-order logic we introduce:\\
\begin{itemize}
    
    \item Variables that represent individual objects.
    \item Functions $\mathcal{F}$ that refer to objects.
    \item Predicates $\mathcal{P}$ - symbols that semantically represents some property or relation.
    \item Quantifiers: universal $\forall$ and existential $\exists$.
\end{itemize}


The Backus Naur form (BNF) for predicate logic:

\[ \phi ::= P(t_1,t_2,...t_n)\ |\ (\neg \phi)\ |\ (\phi \lor \phi)\  |\ (\phi \land \phi)\   |\ (\phi \rightarrow \phi)\ |\ (\forall x \phi)\ |\ (\exists x \phi)\ \]

Where,\\ 
t - called Terms and can be variable, constants or functions on terms  $t ::= x\ |\ c\ |\ f(t,...t)$ 

Atomic formulas are predicates applied to terms: $P(t_1,t_2,...t_n) $

Variables can be free (supplied from outside the formula) or bounded by some quantifier.

Unlike propositional logic, in predicate logical we cannot just evaluate formulas based on given valuation by assigning True or False to the variables.  
We require a model of all function and predicate symbols involved.
A model M for a first‑order logic contains: 
\begin{itemize}
    \item Domain A is a non‑empty set of individuals, mapping each variable to a value in A
    \item Set of functions where n-ary function is interpreted as a concrete function $f^m : A^n \rightarrow A$.
    \item Set of predicates n-ary predicate is interpreted as a relation $P^m \subseteq A^n $ 
\end{itemize}

Truth in a model is defined recursively:

\begin{itemize}
    \item $ M \models P(t_1,...,t_n)$ iff $(t_1,...,t_n) \in P^m $
    \item $M\models\forall x\phi $ iff for every a $ \in A, M \models [x \rightarrow a ] \phi $
    \item $M\models\exists x\phi $ iff for some a $ \in A, M \models [x \rightarrow a ] \phi $
\end{itemize}

A sentence has no free variables, so its truth value depends solely on M.  

We write $\Gamma\models\psi$ (semantic entailment) when every model that makes all formulas in $\Gamma$ true also makes $\psi $ true. 

Soundness and completeness ensure $\Gamma\vdash\psi$ (provability) exactly when $\Gamma\models\psi$,
but the two perspectives complement each other: proofs are constructive, while counter‑models establish non‑entailment.



\subsection {Model checking}

Model checking is a core approach in Formal verification that involves an algorithmic search of the state space of a system model to check if a
given property (specification) holds, where:\\
Model - is an abstract representation of the system (often a state-transition graph or a
program) capturing its possible states and transitions. The model can be described in a modeling language
(like Promela for SPIN, or as a transition system in a tool like NuSMV). \\
Specification - is a formal property that the model should satisfy, commonly expressed in a temporal logic such as LTL or CTL. For example, an LTL
specification might state that “whenever request happens, eventually response occurs” or a CTL property
might require that “in all execution paths, a certain error state is never reached.” \\ 
\subsection {Classification of Model checking methods}
\subsection {Comparison of Model checking methods}

State explosion - the exponential growth of the state space with system parameters (like number of components or data bits).
Combating state explosion is a central theme in model checking research


\subsection {Application of Model checking}
\subsubsection{Hardware model checking}
\subsubsection{Software model checking}
\subsubsection{Reactive systems model checking}

\section {Self-replicating robotic systems}

Reactive system: a system that maintains an
ongoing interaction with its environment, responding to inputs with outputs continually (examples include
operating systems, web servers, control systems). Reactive systems are typically non-terminating and
concurrency-intensive. 
